\documentclass[12p]{article}

\begin{document}
\title{FLIP: Pr\'esentation de transparents par IP Multicast\\(Technologie Push)}
\author{Responsable: J\"org Nonnenmacher\\Etudiants: Denis Arnaud (Casier 6) et S\'ebastien Fabre (Casier 30)}
\maketitle

\section{Introduction}
Une nouvelle technologie appara\^{\i}t prometteuse pour Internet. Elle est
d\'enomm\'ee {\it Push} et correspond \`a l'utilisation de son terminal
comme d'un r\'ecepteur, en continu, d'un flux d'information provenant de
diff\'erents serveurs. On ne va plus chercher l'information sur le serveur:
cette derni\`ere est directement d\'elivr\'ee sur la console du
r\'ecepteur. Cette technologie peut par exemple permettre de recevoir chez
soi les cours de la bourse, des d\'ep\^eches AFP, ou encore le bulletin
m\'et\'eo du lendemain.

Actuellement, l'on trouve sur Internet des techniques permettant ce type
d'applications, mais cela oblige les navigateurs \`a aller chercher, en
permanence, l'information sur les serveurs, et \`a la rapatrier. Il s'agit
donc en fait, \`a proprement parler, de {\it Pull} et non de {\it Push}.

L'on se propose, dans ce projet, d'utiliser le r\'eseau Multicast afin de
diffuser un flux d'information (en l'occurence, ici, une pr\'esentation de
transparents) d'un \'emetteur vers plusieurs r\'ecepteurs. Et cette
mani\`ere de proc\'eder correspond donc bien \`a la technologie {\it Push}.

\section{Description}
L'\'emetteur poss\`ede le contr\^ole du flux d'information, c'est-\`a-dire
qu'il s'abonne \`a un canal Multicast. L'\'emetteur choisit les
transparents et les diffuse sur le canal IP Multicast dans l'ordre qui lui
convient, avec le rythme qu'il d\'esire (possibilit\'e d'aller en avant ou
en arri\`ere, et de rendre la s\'equence automatique).

Le r\'ecepteur, lui, ne peut que s'abonner \`a un canal Multicast, sur
lequel il recevra le flux d'information. Il ne peut pas exercer de
contr\^ole sur la r\'eception des transparents.

Les transparents se pr\'esentent sous la forme de pages HTML
classiques. Aussi, \'emetteur comme r\'ecepteur doivent disposer d'un
navigateur pour les visualiser. Le transfert sur le canal Multicast, par
contre, s'effectue sous la forme de blocs d'octets. Il y a donc conversion
des donn\'ees, entre leur utilisation et leur transport.

Le projet sera r\'ealis\'e dans le langage de programmation Java. De plus,
il devra comporter plusieurs parties, correspondant aux deux couches
logiques suivantes:
\begin{itemize}
  \item une couche de contr\^ole (d\'efilement de la s\'equence des
    transparents, et abonnement \`a Multicast),
  \item une couche de conversion des donn\'ees (HTML $\rightarrow$ Blocs
    d'octets, et Blocs d'octet $\rightarrow$ HTML) pour les rendre
    transportables sur Multicast,
  \item une interface graphique:
    \begin{itemize}
      \item Pour l'\'emetteur: panel de contr\^ole de d\'efilement, gestion
        du nom de canal Multicast,
      \item pour le r\'ecepteur: afficher les transparents, les enregistrer
        et les repasser ult\'erieurement.
    \end{itemize}
\end{itemize}
  
\section{Cheminement}
La partie correspondant au transport de blocs d'octets sur le canal
Multicast a d\'ej\`a \'et\'e programm\'ee.

Aussi, dans un premier temps, on essaiera d'effectuer le simple transfert
de pages HTML texte, puis d'inclure dans ces derni\`eres des images et,
\'eventuellement, des applets Java. Ce module n\'ecessite donc de
programmer la conversion des donn\'ees HTML/blocs d'octets.

Puis, ou parall\`element, on d\'eveloppera l'interface graphique et la
partie de contr\^ole de diffusion.

Le travail sera d\'ecompos\'e en modules, s'articulant les uns par rapport
aux autres gr\^ace \`a des interfaces claires. La programmation des modules
se verra donc r\'epartie entre les deux \'el\`eves du bin\^ome.
\end{document}
